\chapter{Le gestionnaire d'actifs et le fonds}\label{ch:fm:the-asset-manager-and-the-fund}

En date du 3 juin 2019, un fonds britannique à revenus en actions a suspendu ses opérations. Deux ans plus tôt, il valait plus de 10.1~milliards de livres, et il valait 3.6~milliards
de livres à la veille de la suspension~; ses investisseurs n'avaient cessé de réclamer leur argent, et ses gérants n'avaient cessé de vendre ce qui était facile à vendre.
Le régulateur britannique a constaté par la suite qu'au moment de la suspension, seuls 8~\% des investissements du fonds pouvaient être vendus dans un délai de sept jours, alors que ses investisseurs
avaient droit à leur argent dans un délai de quatre jours. Les conditions de liquidité d'un fonds sont une promesse portant sur ses actifs, consignée dans les documents d'un véhicule~; ce
chapitre lit le véhicule, ses frais et ses promesses tels que le gestionnaire doit les faire fonctionner.

\section{Structures~: gestionnaire, fonds et véhicule}

Une société de gestion (livre~1, chapitre~1) ne possède pas l'argent qu'elle gère. Elle fait fonctionner un fonds, un véhicule juridique distinct appartenant à ses investisseurs,
dans le cadre d'un mandat (livre~1, chapitre~3), et se rémunère par des frais prélevés sur ce fonds. La séparation protège les investisseurs si la société de gestion fait défaut, et permet
de proposer une même stratégie sous plusieurs enveloppes.

\begin{definition}[Commandité, commanditaire]\label{def:fm:the-asset-manager-and-the-fund:gp}
Beaucoup de fonds privés sont des sociétés en commandite. Le \emph{commandité}\index{commandité} (general partner) gère la société en commandite, répond de ses dettes
et délègue en général les décisions d'investissement au gestionnaire, qui lui est affilié. Les \emph{commanditaires}\index{commanditaire} (limited partners) sont les investisseurs~:
ils apportent du capital, partagent les gains et les pertes, et ne répondent qu'à hauteur de ce qu'ils ont apporté.
\end{definition}

\begin{definition}[Structure maître--nourricier]\label{def:fm:the-asset-manager-and-the-fund:master}
Une \emph{structure maître--nourricier}\index{structure maître--nourricier} (master--feeder) est un fonds dans lequel les investisseurs souscrivent à un ou plusieurs fonds nourriciers, chacun
conçu pour un groupe d'investisseurs (investisseurs imposables résidents, investisseurs offshore ou exonérés d'impôt), et les nourriciers investissent l'intégralité de leurs actifs dans un seul
fonds maître, qui détient le portefeuille.
\end{definition}

\begin{definition}[OPCVM, compte géré séparément]\label{def:fm:the-asset-manager-and-the-fund:ucits}
Un \emph{OPCVM}\index{OPCVM} (UCITS) est un fonds agréé au titre de la directive de l'Union européenne sur les organismes de placement collectif en
valeurs mobilières~: il peut être vendu au public dans toute l'Union, et en contrepartie il doit investir dans des actifs éligibles et
dans les limites de diversification, et racheter ses parts à la demande de tout investisseur, en ne suspendant qu'exceptionnellement ses rachats. Un \emph{compte géré
séparément}\index{compte géré séparément} (managed account) est un portefeuille détenu au nom propre d'un seul investisseur et géré
dans le cadre d'un contrat de gestion, en dehors de tout fonds mutualisé.
\end{definition}

\begin{omfigure}
\begin{tikzpicture}[font=\footnotesize,>=stealth,box/.style={draw,rounded corners,minimum height=0.62cm,align=center,fill=omDef!10},
  svc/.style={draw,rounded corners,minimum height=0.62cm,align=center,fill=gray!12}]
\node[box,minimum width=2.8cm] (on) at (-2.2,3.4) {nourricier onshore\\{\scriptsize investisseurs imposables}};
\node[box,minimum width=2.8cm] (off) at (2.2,3.4) {nourricier offshore\\{\scriptsize non-résidents, exonérés d'impôt}};
\node[box,minimum width=3.0cm,fill=omProp!12] (m) at (0,1.8) {fonds maître\\{\scriptsize détient le portefeuille}};
\node[box,minimum width=2.4cm,fill=omThm!10] (mgr) at (-4.6,1.8) {gestionnaire\\{\scriptsize décisions d'investissement}};
\node[svc,minimum width=2.2cm] (adm) at (-3.3,0) {administrateur\\{\scriptsize VL, registre}};
\node[svc,minimum width=2.2cm] (cus) at (-0.9,0) {dépositaire\\{\scriptsize conservation}};
\node[svc,minimum width=2.2cm] (pb) at (1.5,0) {prime broker\\{\scriptsize financement}};
\node[svc,minimum width=2.2cm] (aud) at (4.0,0) {commissaire\\aux comptes\\{\scriptsize comptes annuels}};
\draw[->] (on) -- (m); \draw[->] (off) -- (m);
\draw[->,omThm] (mgr) -- node[above,font=\scriptsize] {mandat} (m);
\foreach \s in {adm,cus,pb,aud} {\draw[<->,gray] (m.south) -- (\s.north);}
\end{tikzpicture}

\omcaption{Un hedge fund en \omterm{def:fm:the-asset-manager-and-the-fund:master}{structure maître--nourricier}. Les investisseurs souscrivent au nourricier qui leur est destiné~; les nourriciers investissent dans le fonds maître, qui détient le
portefeuille. Le gestionnaire prend les décisions~; l'administrateur calcule la valeur liquidative et tient le registre, le \omterm{def:fm:the-asset-manager-and-the-fund:admin}{dépositaire} conserve les actifs, le prime broker
finance les positions, et le commissaire aux comptes certifie les comptes.}\label{fig:fm:the-asset-manager-and-the-fund:structure}
\end{omfigure}

\section{Prestataires de services~: administrateur de fonds, dépositaire, commissaire aux comptes, prime broker}

\begin{definition}[Administrateur de fonds, dépositaire]\label{def:fm:the-asset-manager-and-the-fund:admin}
Un \emph{administrateur de fonds}\index{administrateur de fonds} (fund administrator) est la société, indépendante du gestionnaire, qui calcule la valeur liquidative d'un fonds, tient
son registre des investisseurs, traite les souscriptions et les rachats et calcule les frais de chaque investisseur. Un \emph{dépositaire}\index{dépositaire} (custodian) est
l'établissement qui conserve les actifs d'un fonds, dans des comptes séparés de ses propres comptes et de ceux du gestionnaire.
\end{definition}

Les prestataires de services sont les contrôles du fonds. Un administrateur qui valorise le portefeuille de manière indépendante empêche un gestionnaire de valoriser lui-même son propre livre~; un
\omterm{def:fm:the-asset-manager-and-the-fund:admin}{dépositaire} qui détient les actifs empêche qu'ils servent à autre chose qu'au fonds~; un commissaire aux comptes vérifie les deux une fois par an. Un prime broker (livre~1, chapitre~6)
finance les positions et détient souvent les actifs qu'il prête en garantie, ce qui explique que les documents d'un hedge fund précisent ce qui peut être
réutilisé (rehypothecation, chapitre~18). Le gestionnaire choisit ces sociétés, mais la diligence raisonnable des investisseurs (chapitre~25) contrôle chacune d'elles.

\section{Les frais en pratique~: catégories de parts, seuils de déclenchement, cristallisation et égalisation}

Les frais d'un fonds sont les frais de gestion et les commissions de performance du livre~1, chapitre~3, assortis d'un high-water mark~; en pratique, ils sont fixés par \omterm{def:fm:the-asset-manager-and-the-fund:class}{catégorie de parts},
et trois mécanismes déterminent qui paie quoi.

\begin{definition}[Catégorie de parts, seuil de déclenchement des commissions, cristallisation]\label{def:fm:the-asset-manager-and-the-fund:class}
Une \emph{catégorie de parts}\index{catégorie de parts} (share class) est un ensemble de parts d'un fonds, avec ses propres conditions~: devise, frais, liquidité, investissement minimum. Un
\emph{seuil de déclenchement des commissions}\index{seuil de déclenchement des commissions} (hurdle) est un rendement, fixe ou lié à un indice de référence ou à un taux monétaire, que le fonds doit dépasser avant qu'une commission de performance
soit due. La \emph{cristallisation}\index{cristallisation} est le moment où une commission de performance provisionnée devient exigible pour le gestionnaire et ne peut plus
être annulée par des pertes ultérieures~; les documents du fonds en fixent la fréquence.
\end{definition}

Plus les commissions cristallisent souvent, plus le gestionnaire gagne sur les mêmes rendements~: un gain cristallisé en mars n'est pas restitué lorsque avril est en perte.
Sur cinq ans de rendements mensuels, avec une moyenne de 8~\% par an et une volatilité de 12~\%, une commission de performance de 20~\% cristallisée chaque mois prélève en moyenne 11.54~\%
du capital de départ sur 1\,000 trajectoires, contre 11.26~\% en \omterm{def:fm:the-asset-manager-and-the-fund:class}{cristallisation} trimestrielle et 10.78~\% en \omterm{def:fm:the-asset-manager-and-the-fund:class}{cristallisation} annuelle.

\begin{definition}[Égalisation]\label{def:fm:the-asset-manager-and-the-fund:eq}
L'\emph{égalisation}\index{egalisation@égalisation} désigne l'ensemble des ajustements, ou la pratique consistant à émettre une série distincte de parts pour chaque date de souscription,
grâce auxquels chaque investisseur d'un fonds mutualisé paie des commissions de performance sur les gains de son propre investissement, mesurés à partir de son propre prix d'entrée.
\end{definition}

\begin{example}[Le passager clandestin]\label{ex:fm:the-asset-manager-and-the-fund:free}
L'investisseur A souscrit 100 à une valeur liquidative de 1.00. Le fonds baisse de 10~\%~; l'investisseur B souscrit 100 à 0.90. Le fonds remonte ensuite à 1.05. Avec
un seul prix de part mutualisé et un high-water mark à 1.00, le fonds prélève 20~\% des 0.05 par part au-dessus du seuil~: A paie 1.00, B paie 1.11. Le
gain propre de B va de 0.90 à 1.05~; avec une comptabilité par séries, B paie 3.33. Sans \omterm{def:fm:the-asset-manager-and-the-fund:eq}{égalisation}, B profite gratuitement de 2.22 de commission de performance que la perte de A
a engendrée. Le cas inverse, un investisseur qui souscrit au-dessus du seuil et paie sur des gains qui ne sont pas les siens, est tout aussi injuste et tout aussi fréquent.
\end{example}

\begin{method}[Vérifier les conditions de frais d'un fonds]\label{met:fm:the-asset-manager-and-the-fund:fees}
\begin{enumerate}
\item Lister chaque \omterm{def:fm:the-asset-manager-and-the-fund:class}{catégorie de parts}, avec ses frais de gestion, sa commission de performance, son seuil de déclenchement, son high-water mark et la fréquence de \omterm{def:fm:the-asset-manager-and-the-fund:class}{cristallisation}.
\item Vérifier si les commissions de performance sont calculées par investisseur (séries ou \omterm{def:fm:the-asset-manager-and-the-fund:eq}{égalisation}) ou sur le prix mutualisé de la part.
\item Simuler le prélèvement des frais sur des trajectoires de rendements de la stratégie, et non sur son rendement espéré~: la \omterm{def:fm:the-asset-manager-and-the-fund:class}{cristallisation} et le high-water mark agissent sur la trajectoire.
\item Comparer le résultat avec l'exemple chiffré de la note d'information~; l'administrateur calcule les frais réels et doit le reproduire.
\end{enumerate}
\end{method}

\section{Conditions de liquidité et inadéquation de liquidité}

\begin{definition}[Préavis de rachat, période de blocage]\label{def:fm:the-asset-manager-and-the-fund:notice}
Le \emph{préavis de rachat}\index{préavis de rachat} (redemption notice period) d'un fonds est le délai que doit respecter un investisseur entre la demande de remboursement de ses parts et
la date de traitement à laquelle elles sont rachetées. La \emph{période de blocage}\index{période de blocage} (lock-up) est une période initiale suivant une souscription, pendant laquelle
les parts ne peuvent pas être rachetées, ou seulement moyennant une pénalité.
\end{definition}

\begin{definition}[Plafonnement des rachats, side pocket, swing pricing]\label{def:fm:the-asset-manager-and-the-fund:gate}
Un \emph{plafonnement des rachats}\index{plafonnement des rachats} (redemption gate) limite la part des actifs d'un fonds (ou d'un investisseur) qui peut être rachetée à une même date de traitement~;
les demandes au-delà de ce seuil sont reportées. Un \emph{side pocket}\index{side pocket} est un compte séparé, ou une \omterm{def:fm:the-asset-manager-and-the-fund:class}{catégorie de parts}, dans lequel sont placés des actifs peu liquides ou
difficiles à valoriser, rachetables seulement lorsqu'ils sont vendus, afin que les investisseurs qui entrent et sortent ne les échangent pas à un prix estimé.
Le \emph{swing pricing}\index{swing pricing} ajuste le prix auquel les investisseurs souscrivent ou rachètent les jours de flux nets importants, de sorte que le
coût de transaction qu'ils occasionnent soit supporté par eux et non par les investisseurs qui restent.
\end{definition}

\begin{definition}[Inadéquation de liquidité]\label{def:fm:the-asset-manager-and-the-fund:mismatch}
Un fonds présente une \emph{inadéquation de liquidité}\index{inadéquation de liquidité} (liquidity mismatch) lorsque ses investisseurs peuvent retirer, sur une période donnée, une part de sa valeur supérieure
à celle qu'il pourrait vendre sur cette même période à un coût normal.
\end{definition}

Le Form PF, rapport confidentiel que déposent les grands conseillers américains de hedge funds, donne la vue d'ensemble~: les investisseurs des hedge funds éligibles pouvaient
retirer 9.0~\% de leur capital en sept jours au troisième trimestre 2025, alors que les fonds déclaraient pouvoir vendre 51.8~\% de leurs
portefeuilles en sept jours (voir la \cref{fig:fm:the-asset-manager-and-the-fund:liquidity}). À chaque horizon, les portefeuilles étaient plus liquides que les
conditions offertes aux investisseurs~: globalement, l'industrie des hedge funds n'est pas en inadéquation. Le fonds suspendu en 2019 présentait une inadéquation de plus de douze fois
à sept jours~: ses investisseurs auraient pu récupérer l'intégralité de leurs avoirs en quatre jours, alors qu'il ne pouvait vendre que 8~\% de ses actifs en sept jours.

\begin{omfigure}
\begin{tikzpicture}
\begin{axis}[width=10.5cm,height=5.6cm,xmode=log,log basis x=10,xmin=0.8,xmax=450,ymin=0,ymax=105,
  xtick={1,7,30,90,180,365},xticklabels={1,7,30,90,180,365},xlabel={\small jours (échelle logarithmique)},ylabel={\small part cumulée de la valeur liquidative (\%)},
  tick label style={font=\footnotesize},grid=major,grid style={gray!20},table/col sep=comma,
  legend style={font=\scriptsize,at={(0.5,-0.3)},anchor=north,/tikz/every even column/.append style={column sep=8pt}},legend columns=2]
\addplot[omThm,thick,mark=*] table[x=days,y=investor]{figdata/desk/04-the-asset-manager-and-the-fund/liquidity.csv};
\addplot[omDef,thick,mark=square*] table[x=days,y=portfolio]{figdata/desk/04-the-asset-manager-and-the-fund/liquidity.csv};
\addplot[gray,thick,dashed,mark=triangle*] table[x=days,y=illiquid]{figdata/desk/04-the-asset-manager-and-the-fund/liquidity.csv};
\addplot[only marks,mark=diamond*,mark size=3pt,omProp] coordinates {(4,100)};
\legend{les investisseurs des hedge funds peuvent retirer,les portefeuilles des hedge funds peuvent être vendus,{fonds peu liquide (8~\% en 7 jours)},{ses investisseurs~: tous dans les 4 jours}}
\end{axis}
\end{tikzpicture}

\omcaption{Liquidité des conditions offertes aux investisseurs comparée à celle du portefeuille. Hedge funds américains éligibles, agrégés, T3 2025 (Form PF,
SEC Private Funds Statistics, tableaux 8.22 et 8.23)~; et un fonds illustratif dont 8~\% de ses actifs peuvent être vendus en sept jours, soit le chiffre
que le régulateur britannique a donné pour le fonds suspendu en 2019, dont les investisseurs pouvaient racheter l'intégralité de leurs avoirs en quatre jours.}\label{fig:fm:the-asset-manager-and-the-fund:liquidity}
\end{omfigure}

\begin{dated}{2026-09}{Le registre public de la liquidité des fonds}\label{dat:fm:the-asset-manager-and-the-fund:pf}
SEC Private Funds Statistics, troisième trimestre 2025 (hedge funds éligibles, valeur liquidative agrégée de 4\,750~milliards de dollars)~: actifs susceptibles d'être suspendus, 3\,182
milliards de dollars, et susceptibles d'être plafonnés, 2\,053~milliards de dollars~; actifs plafonnés, 121~milliards de dollars, placés en \omterm{def:fm:the-asset-manager-and-the-fund:gate}{side pocket}, 103~milliards de dollars, et suspendus, 25~milliards de dollars (2.5~\%, 2.2~\% et
0.5~\% de la valeur liquidative). UCITS Directive 2009/65/EC, article 84~: un \omterm{def:fm:the-asset-manager-and-the-fund:ucits}{OPCVM} rachète à la demande de tout porteur de parts et ne peut suspendre que temporairement, dans
des cas exceptionnels et dans l'intérêt des porteurs. Fonds ouverts américains, SEC Rule 22e-4~: un investissement peu liquide est un actif qui ne peut pas être vendu en
sept jours calendaires sans modifier sensiblement son cours, et un fonds ne peut en acquérir un si plus de 15~\% de ses actifs nets deviendraient alors peu liquides.
La Rule 22c-1(a)(3) autorise le \omterm{def:fm:the-asset-manager-and-the-fund:gate}{swing pricing} au-delà d'un seuil de flux nets.
\end{dated}

\begin{dated}{2026-09}{La suspension de 2019 dans les mots du régulateur}\label{dat:fm:the-asset-manager-and-the-fund:weif}
Financial Conduct Authority, 11 avril 2024~: le directeur général agréé du fonds «~n'a pas su gérer la liquidité du fonds~» entre
31 juillet 2018 et la suspension du 3 juin 2019~; les investisseurs présents dans le fonds à la date de la suspension perçoivent une part d'un dispositif d'indemnisation pouvant atteindre
230~millions de livres, approuvé par la High Court en février 2024. FCA, 5 août 2025~: le régulateur a décidé d'infliger une amende de 40~millions de livres à la société de gestion et une amende de
5\,888\,800~livres à son fondateur, et de lui interdire de gérer des fonds de détail~; les deux parties ont porté les décisions devant l'Upper Tribunal, de sorte que les conclusions restent provisoires. Le
régulateur constate qu'à partir de juillet 2018 le fonds «~a vendu de manière disproportionnée les placements les plus liquides~» et a acheté des actifs moins liquides.
\end{dated}

\begin{proposition}[Vendre d'abord les actifs liquides laisse les investisseurs restants dans un fonds peu liquide]\label{prop:fm:the-asset-manager-and-the-fund:liquidfirst}
Honorer un rachat portant sur une part $q$ de la valeur liquidative en vendant d'abord les actifs les plus liquides. Si une part $\ell_h$ de la valeur liquidative peut être vendue dans l'horizon $h$ et
$q\le\ell_h$, les investisseurs restants détiennent un fonds dont seule une part $(\ell_h-q)/(1-q)$ peut être vendue dans l'horizon $h$~; si $q\ge\ell_h$, aucune.
\end{proposition}

\begin{proof}
La vente retire $q$ de la poche liquide et rien des autres~; la valeur liquidative restante vaut $1-q$ avant coûts.
\end{proof}

Vendre d'abord les actifs liquides minimise le coût du rachat du jour et laisse aux investisseurs restants un fonds moins liquide~; le rachat suivant devient alors
plus coûteux, et racheter tôt devient plus avantageux que rester. Le modèle de \texttt{firm.fundterms} (voir le \cref{lst:fm:the-asset-manager-and-the-fund:redeem})
répond à un rachat de 30~\% avec un préavis de sept jours, pour deux échelles de liquidité, le portefeuille agrégé des hedge funds du Form PF et le fonds peu liquide, avec
des coûts de vente illustratifs qui doublent pour les actifs vendus plus vite que leur horizon~:

\begin{center}
\small
\begin{tabular}{@{}>{\raggedright\arraybackslash}p{3.0cm}>{\raggedright\arraybackslash}p{2.4cm}>{\centering\arraybackslash}p{2.0cm}>{\centering\arraybackslash}p{3.2cm}>{\centering\arraybackslash}p{3.2cm}@{}}
\toprule
fonds & politique & racheté & coût supporté par les investisseurs restants & vendable dans les 7 jours suivants\\
\midrule
échelle hedge funds & liquide d'abord & 30~\% & 0.04~\% & 31.1~\%\\
 & au prorata & 30~\% & 5.39~\% & 51.8~\%\\
 & plafonnement à 10~\% & 10~\% & 0.01~\% & 46.4~\%\\
 & liquide d'abord, \omterm{def:fm:the-asset-manager-and-the-fund:gate}{swing pricing} & 30~\% & 0 & 31.1~\%\\
fonds peu liquide & liquide d'abord & 30~\% & 1.69~\% & 0\\
 & au prorata & 30~\% & 8.81~\% & 8.0~\%\\
 & plafonnement à 10~\% & 10~\% & 0.09~\% & 0\\
 & liquide d'abord, \omterm{def:fm:the-asset-manager-and-the-fund:gate}{swing pricing} & 30~\% & 0 & 0\\
\bottomrule
\end{tabular}
\end{center}

\omcode{code/firm/fundterms/firm_fundterms.py}{85}{112}{Répondre à un rachat~: liquide d'abord ou une tranche de chaque poche, un plafonnement et swing pricing.}\label{lst:fm:the-asset-manager-and-the-fund:redeem}

\begin{omfigure}
\begin{tikzpicture}
\begin{axis}[width=11cm,height=5.4cm,ybar,bar width=10pt,ymin=0,ymax=60,ylabel={\small \%},
  xtick={0,1,2,3},xticklabels={liquide d'abord,au prorata,plafonnement 10~\%,{liquide d'abord~; swing}},tick label style={font=\footnotesize},xticklabel style={text width=2.3cm,align=center},
  table/col sep=comma,enlarge x limits=0.15,grid=major,grid style={gray!20},
  legend cell align=left,legend style={font=\scriptsize,at={(0.5,-0.25)},anchor=north,/tikz/every even column/.append style={column sep=8pt}},legend columns=2]
\addplot[fill=omDef!55,draw=omDef] table[x=k,y=hf_liquid]{figdata/desk/04-the-asset-manager-and-the-fund/stress.csv};
\addplot[fill=gray!40,draw=gray] table[x=k,y=ill_liquid]{figdata/desk/04-the-asset-manager-and-the-fund/stress.csv};
\addplot[fill=omThm!55,draw=omThm] table[x=k,y=hf_dilution]{figdata/desk/04-the-asset-manager-and-the-fund/stress.csv};
\addplot[fill=omProp!55,draw=omProp] table[x=k,y=ill_dilution]{figdata/desk/04-the-asset-manager-and-the-fund/stress.csv};
\legend{échelle hedge funds~: vendable dans les 7 jours suivants,fonds peu liquide~: vendable dans les 7 jours suivants,échelle hedge funds~: coût pour les investisseurs restants,fonds peu liquide~: coût pour les investisseurs restants}
\end{axis}
\end{tikzpicture}

\omcaption{Rachat de 30~\% avec préavis de sept jours, selon quatre politiques~: part du fonds restant vendable dans les sept jours
suivants, et le coût de vente supporté par les investisseurs restants, en pourcentage de leur valeur liquidative. Échelles~: l'agrégat du Form PF et le fonds
peu liquide illustratif~; coûts de vente illustratifs. Données~: \texttt{fm\_fund.stress\_table}.}\label{fig:fm:the-asset-manager-and-the-fund:stress}
\end{omfigure}

La tranche au prorata conserve la structure du fonds mais coûte aux investisseurs restants 5.4~\% sur l'échelle des hedge funds et 8.8~\% sur celle du fonds peu liquide, car une tranche de
chaque poche doit être vendue dans la semaine. Le \omterm{def:fm:the-asset-manager-and-the-fund:gate}{swing pricing} reporte le coût sur les investisseurs qui rachètent et laisse le problème de liquidité en l'état. Un \omterm{def:fm:the-asset-manager-and-the-fund:gate}{plafonnement
des rachats} gagne du temps~: il verse 10~\% maintenant, reporte 20~\%, et laisse au gérant plusieurs semaines pour vendre les actifs les plus lents à un coût normal. Le fonds peu liquide n'a pas de bonne
option~: quoi qu'il fasse, ses premiers 8~\% de rachats épuisent tout ce qu'il peut vendre en une semaine. Le plus grand rachat qu'il peut honorer en
sept jours à un coût inférieur à 1~\% est de 8~\% de sa valeur liquidative~; celui de l'échelle du Form PF est de 51.8~\%.

\section{Ce que demandent les investisseurs~: lettres d'accompagnement et comptes gérés séparément}

\begin{definition}[Lettre d'accompagnement, clause de la nation la plus favorisée]\label{def:fm:the-asset-manager-and-the-fund:side}
Une \emph{lettre d'accompagnement}\index{lettre d'accompagnement} (side letter) est un accord conclu entre un fonds (ou son gérant) et un seul investisseur qui modifie les conditions du fonds pour
cet investisseur~: frais réduits, information renforcée, meilleure liquidité, engagement de capacité. Une
\emph{clause de la nation la plus favorisée}\index{clause de la nation la plus favorisée} est une promesse selon laquelle l'investisseur se verra offrir toutes les conditions plus avantageuses accordées à un autre investisseur de taille identique ou inférieure.
\end{definition}

Une \omterm{def:fm:the-asset-manager-and-the-fund:side}{lettre d'accompagnement} qui accorde à un investisseur une meilleure liquidité crée une \omterm{def:fm:the-asset-manager-and-the-fund:mismatch}{inadéquation de liquidité} entre investisseurs~: il rachète en premier et laisse les autres avec un fonds moins
liquide, ce qui explique que les régulateurs et les consultants demandent si un investisseur bénéficie d'une liquidité préférentielle. Les gros investisseurs qui veulent une transparence totale, leur propre liquidité et aucune exposition aux flux des autres investisseurs
recourent plutôt à un \omterm{def:fm:the-asset-manager-and-the-fund:ucits}{compte géré séparément}~: le gérant opère
sur le compte propre de l'investisseur selon le même modèle, et le \omterm{def:fm:the-asset-manager-and-the-fund:admin}{dépositaire} de l'investisseur détient les actifs. Le gérant fait alors tourner plusieurs copies d'une même stratégie, et
doit répartir les transactions entre elles de manière équitable, ce que sa fonction de conformité vérifie (chapitre~16).

\begin{remark}[La liquidité est une propriété du couple]\label{rem:fm:the-asset-manager-and-the-fund:pair}
Un fonds n'est pas liquide ou illiquide~; ses actifs et ses conditions sont liquides ou non l'un par rapport à l'autre. Le même portefeuille est sûr dans un véhicule
à négociation trimestrielle avec un préavis de 90~jours, et fragile dans un véhicule à négociation quotidienne. Un gérant qui modifie la liquidité du portefeuille, comme le régulateur
l'a constaté en 2019, modifie la promesse du fonds sans modifier ses documents.
\end{remark}

\section{Tutoriel~: un fonds, deux investisseurs et une mauvaise semaine}\label{tut:fm:the-asset-manager-and-the-fund:tut}

\begin{tutorial}
\textbf{Objectif.} Calculer qui paie quelles commissions, et ce que fait un rachat important aux investisseurs qui restent. \textbf{État final~:} le tableau de stress et
la \cref{fig:fm:the-asset-manager-and-the-fund:stress}.
\begin{tutsteps}
\item \textbf{Commissions par investisseur.} \texttt{fm\_fund.free\_ride()} exécute l'\cref{ex:fm:the-asset-manager-and-the-fund:free} à travers
\texttt{firm.fundterms.pooled\_vs\_series}, qui encapsule les \texttt{firm.fees} du livre~1.
\item \textbf{\omterm{def:fm:the-asset-manager-and-the-fund:class}{Cristallisation}.} \texttt{fm\_fund.crystallisation()} calcule la moyenne des commissions perçues sur 1\,000 trajectoires pour une
\omterm{def:fm:the-asset-manager-and-the-fund:class}{cristallisation} mensuelle, trimestrielle et annuelle.
\item \textbf{Échelles.} \texttt{hedge\_fund\_ladder()} transforme la liquidité de portefeuille cumulée du Form PF en tranches~; \texttt{illiquid\_ladder()}
construit le fonds à 8~\% en sept jours.
\item \textbf{Stress.} \texttt{stress\_table()} honore une demande de 30~\% selon quatre politiques (la \cref{prop:fm:the-asset-manager-and-the-fund:liquidfirst}).
\end{tutsteps}
\end{tutorial}

\textbf{Pour aller plus loin.} Lancer deux semaines de rachats consécutives et observer la deuxième semaine du fonds à liquidité d'abord~; accorder à un investisseur une
\omterm{def:fm:the-asset-manager-and-the-fund:side}{lettre d'accompagnement} prévoyant une liquidité quotidienne dans un fonds à négociation mensuelle.

\section{Construction~: les conditions du fonds}\label{bld:fm:the-asset-manager-and-the-fund:fundterms}

\begin{build}
\bfield{Rôle} Le fonds tel que le voient ses investisseurs et son administrateur~: commissions par investisseur, \omterm{def:fm:the-asset-manager-and-the-fund:class}{cristallisation}, et ce que deviennent ses conditions de liquidité
en cas de rachats.
\bfield{Interface} \texttt{firm.fundterms}~: \texttt{pooled\_vs\_series(returns, entries, terms)}~; \texttt{crystallised(returns, perf, every)}~;
\texttt{Bucket(days, share, cost)}, \texttt{ladder\_share\_within}~; \texttt{redeem(ladder, request, notice\_days, policy, gate, swing, fire)}~;
\texttt{max\_redemption(ladder, notice\_days, max\_cost)}.
\bfield{Règles} L'arithmétique des commissions est celle des \texttt{firm.fees} du livre~1, sans modification~; sans \omterm{def:fm:the-asset-manager-and-the-fund:gate}{swing pricing}, les investisseurs restants supportent tout le coût de vente~;
un plafonnement reporte les rachats, il ne les annule jamais.
\bfield{Tests d'acceptation} \texttt{code/firm/fundterms/tests/}~: le free ride apparaît dans le fonds groupé et non dans les séries~; une \omterm{def:fm:the-asset-manager-and-the-fund:class}{cristallisation} plus fréquente
ne réduit jamais la commission~; la liquidité d'abord coûte moins cher et laisse moins de liquidité que le prorata~; le plafonnement reporte le reste~; le \omterm{def:fm:the-asset-manager-and-the-fund:gate}{swing pricing}
supprime la dilution des investisseurs restants.
\bfield{Pour aller plus loin} Plusieurs dates de négociation avec des demandes différées mises en file~; une \omterm{def:fm:the-asset-manager-and-the-fund:gate}{side pocket} qui retire une tranche de la valeur liquidative rachetable~; un investisseur
disposant d'une \omterm{def:fm:the-asset-manager-and-the-fund:side}{lettre d'accompagnement}.
\end{build}

\begin{omsources}
\item Financial Conduct Authority, communiqués de presse des 11 avril 2024 (Link Fund Solutions) et 5 août 2025 (Woodford Equity Income Fund).
\item US Securities and Exchange Commission, Division of Investment Management, \emph{Private Funds Statistics}, troisième trimestre civil 2025.
\item Directive 2009/65/EC (UCITS), article 84~; 17 CFR 270.22e-4 et 270.22c-1 (gestion du risque de liquidité et swing pricing aux États-Unis).
\end{omsources}

\section{Exercices}

\begin{exercise}[$\star$]\label{exo:fm:the-asset-manager-and-the-fund:1}
Dans l'\cref{ex:fm:the-asset-manager-and-the-fund:free}, calculer la commission de performance de chaque investisseur selon le prix de part groupé puis selon la
comptabilité par séries.
\end{exercise}

\begin{exercise}[$\star$]\label{exo:fm:the-asset-manager-and-the-fund:2}
D'après les chiffres du Form PF du T3 2025, quelles parts de la valeur liquidative les investisseurs des hedge funds éligibles pourraient-ils retirer, et leurs portefeuilles vendre, dans les 30~jours~?
\end{exercise}

\begin{exercise}[$\star$]\label{exo:fm:the-asset-manager-and-the-fund:3}
Les hedge funds éligibles ont déclaré 121~milliards de dollars plafonnés, 103~milliards de dollars en \omterm{def:fm:the-asset-manager-and-the-fund:gate}{side pocket} et 25~milliards de dollars suspendus, sur une valeur liquidative de 4\,750~milliards de dollars. Exprimer chacun
de ces montants en part de la valeur liquidative.
\end{exercise}

\begin{exercise}[$\star\star$]\label{exo:fm:the-asset-manager-and-the-fund:4}
Un fonds peut vendre 20~\% de sa valeur liquidative en sept jours et reçoit une demande de rachat de 15~\%, honorée en priorité sur les actifs liquides. Quelle part du fonds restant peut être vendue
dans les sept jours~?
\end{exercise}

\begin{exercise}[$\star\star$]\label{exo:fm:the-asset-manager-and-the-fund:5}
Pourquoi une commission de performance cristallisée chaque mois prélève-t-elle plus qu'une commission cristallisée chaque année sur les mêmes rendements~? Donner un exemple sur deux mois.
\end{exercise}

\begin{exercise}[$\star\star$]\label{exo:fm:the-asset-manager-and-the-fund:6}
Dans la \cref{fig:fm:the-asset-manager-and-the-fund:stress}, pourquoi le prorata coûte-t-il davantage aux investisseurs restants que la liquidité d'abord, tout en leur laissant un
fonds plus liquide~?
\end{exercise}

\begin{exercise}[$\star\star\star$]\label{exo:fm:the-asset-manager-and-the-fund:7}
\emph{Programmation.} Exécuter \texttt{firm.fundterms.redeem} sur l'échelle du fonds peu liquide pour une demande de 5~\%, puis une deuxième demande de 5~\% sur l'échelle restante après la
première, les deux en liquidité d'abord avec un préavis de sept jours. Quelle part peut être vendue dans les sept jours après chacune~?
\end{exercise}

\begin{exercise}[$\star\star\star$]\label{exo:fm:the-asset-manager-and-the-fund:8}
\emph{Trouver la faille.} «~Notre fonds peut vendre 60~\% de ses actifs en une semaine et nos investisseurs disposent d'une liquidité mensuelle avec un préavis de 30~jours, donc nous
n'avons aucun risque de liquidité.~»
\end{exercise}

\section{Problème~: la file du vendredi}

\begin{problem}[{Problème du week-end --- la file du vendredi}]\label{pb:fm:the-asset-manager-and-the-fund:1}
Le vendredi, un fonds ouvert à négociation hebdomadaire avec un préavis de sept jours reçoit des demandes de rachat portant sur 30~\% de sa valeur liquidative. Il faut gérer sa liquidité.

\medskip\noindent\textbf{Partie I --- Le véhicule.}
\begin{enumerate}
\item Définir une \omterm{def:fm:the-asset-manager-and-the-fund:master}{structure maître--nourricier} ainsi que les rôles du \omterm{def:fm:the-asset-manager-and-the-fund:gp}{commandité} et des \omterm{def:fm:the-asset-manager-and-the-fund:gp}{commanditaires}.
\item Que font l'\omterm{def:fm:the-asset-manager-and-the-fund:admin}{administrateur de fonds} et le \omterm{def:fm:the-asset-manager-and-the-fund:admin}{dépositaire}, et pourquoi doivent-ils être indépendants du gérant~?
\item Définir un \omterm{def:fm:the-asset-manager-and-the-fund:notice}{préavis de rachat}, une \omterm{def:fm:the-asset-manager-and-the-fund:notice}{période de blocage}, un \omterm{def:fm:the-asset-manager-and-the-fund:gate}{plafonnement des rachats}, une \omterm{def:fm:the-asset-manager-and-the-fund:gate}{side pocket} et le \omterm{def:fm:the-asset-manager-and-the-fund:gate}{swing pricing}.
\item Qu'impose la directive \omterm{def:fm:the-asset-manager-and-the-fund:ucits}{OPCVM} au sujet des rachats~?
\end{enumerate}

\medskip\noindent\textbf{Partie II --- Le registre public.}
\begin{enumerate}[resume]
\item Donner la liquidité des investisseurs et du portefeuille à sept jours du Form PF pour le T3 2025.
\item Le secteur des hedge funds présente-t-il une inadéquation globale~? À quel horizon l'écart est-il le plus faible en valeur relative~?
\item Donner la part de la valeur liquidative des hedge funds éligibles soumise à \omterm{def:fm:the-asset-manager-and-the-fund:gate}{plafonnement des rachats}, placée en \omterm{def:fm:the-asset-manager-and-the-fund:gate}{side pocket} ou suspendue.
\item Qu'a constaté le régulateur au sujet du fonds suspendu en 2019, et quel est le statut de ses décisions~?
\end{enumerate}

\medskip\noindent\textbf{Partie III --- La file d'attente.}
\begin{enumerate}[resume]
\item Énoncer et démontrer la \cref{prop:fm:the-asset-manager-and-the-fund:liquidfirst}.
\item Pour l'échelle du Form PF, donner le coût pour les investisseurs restants et la part pouvant être vendue dans les sept jours suivants, selon la règle de la liquidité d'abord et selon la règle au prorata.
\item Même question pour le fonds peu liquide.
\item Que reportent, que versent et que laissent liquide un plafonnement de 10~\% pour chaque échelle~?
\item Que change le \omterm{def:fm:the-asset-manager-and-the-fund:gate}{swing pricing}, et que ne change-t-il pas~?
\item Donner le rachat le plus important que chaque échelle peut honorer dans les sept jours, à un coût inférieur à 1~\%.
\end{enumerate}

\medskip\noindent\textbf{Partie IV --- Les commissions.}
\begin{enumerate}[resume]
\item Définir l'\omterm{def:fm:the-asset-manager-and-the-fund:eq}{égalisation} et calculer le passager clandestin dans l'\cref{ex:fm:the-asset-manager-and-the-fund:free}.
\item Donner la commission moyenne perçue pour une \omterm{def:fm:the-asset-manager-and-the-fund:class}{cristallisation} mensuelle, trimestrielle et annuelle.
\item Pourquoi une \omterm{def:fm:the-asset-manager-and-the-fund:side}{lettre d'accompagnement} offrant une meilleure liquidité constitue-t-elle un risque de liquidité pour les autres investisseurs~?
\item Définir une \omterm{def:fm:the-asset-manager-and-the-fund:mismatch}{inadéquation de liquidité}.
\item Énoncer le \emph{résultat nommé}~: le rachat le plus important que chaque fonds peut honorer dans son préavis, à un coût inférieur à 1~\%, ainsi que
la dilution subie par les investisseurs restants sans \omterm{def:fm:the-asset-manager-and-the-fund:gate}{swing pricing} lorsque 30~\% sont honorés en liquidité d'abord.
\item En deux phrases, qu'aurait dû faire le conseil d'administration du fonds peu liquide avant le vendredi~?
\end{enumerate}
\end{problem}

\section{Questions d'entretien}

\begin{interviewq}[$\star$ \iqroles{chercheur}]\label{iq:fm:the-asset-manager-and-the-fund:1}
À quoi sert une \omterm{def:fm:the-asset-manager-and-the-fund:master}{structure maître--nourricier}, et qui détient le portefeuille~?
\end{interviewq}

\begin{interviewq}[$\star$ \iqroles{trader, risque}]\label{iq:fm:the-asset-manager-and-the-fund:2}
Qu'est-ce qu'un \omterm{def:fm:the-asset-manager-and-the-fund:gate}{plafonnement des rachats}, et pourquoi des investisseurs pourraient-ils préférer un fonds qui en dispose~?
\end{interviewq}

\begin{interviewq}[$\star\star$ \iqroles{chercheur}]\label{iq:fm:the-asset-manager-and-the-fund:3}
Un investisseur souscrit dans un fonds mutualisé 10~\% en dessous de son high-water mark. Expliquez le passager clandestin et la manière dont un fonds l'empêche.
\end{interviewq}

\begin{interviewq}[$\star\star$ \iqroles{risque}]\label{iq:fm:the-asset-manager-and-the-fund:4}
Vous devez honorer un rachat important. Pourquoi vendre d'abord les actifs les plus liquides n'est-il pas neutre pour les investisseurs qui restent~?
\end{interviewq}

\begin{interviewq}[$\star\star$ \iqroles{trader}]\label{iq:fm:the-asset-manager-and-the-fund:5}
Que fait le \omterm{def:fm:the-asset-manager-and-the-fund:gate}{swing pricing}, et quel problème ne résout-il pas~?
\end{interviewq}

\begin{interviewq}[$\star\star\star$ \iqroles{risque, chercheur}]\label{iq:fm:the-asset-manager-and-the-fund:6}
Comment mesureriez-vous l'\omterm{def:fm:the-asset-manager-and-the-fund:mismatch}{inadéquation de liquidité} d'un fonds, et quelles données demanderiez-vous au gérant~?
\end{interviewq}
